\subsection{Optimización Específica para Grover mediante Parametrización Reducida}
\label{subsec:grover_reducido_impl}

Además del backend general basado en bloques comprimidos, el desarrollo incorporó una
optimización específica para el algoritmo de Grover. Esta variante parte de la
observación de que, en ciertos escenarios, la evolución del estado puede describirse
mediante una parametrización reducida asociada a la estructura del algoritmo, sin
necesidad de operar constantemente sobre el vector completo de amplitudes, 
su derivación matemática se puede encontrar en el apéndice \ref{ap:grover_reducido}.

Su propósito no fue sustituir la implementación general, sino ofrecer una referencia
complementaria para contrastar el efecto de una simplificación estructural específica
frente al almacenamiento comprimido aplicable de manera general. Mientras el backend
basado en Blosc2 conserva una representación explícita del estado y modifica su forma
de almacenamiento, la variante reducida aprovecha una regularidad particular de Grover
y, por tanto, constituye una optimización especializada.

Esta inclusión resulta útil porque permite distinguir entre dos fuentes distintas de mejora:
por un lado, las obtenidas mediante compresión sin pérdida sobre el vector de estado; por
otro, las derivadas de una representación más compacta sustentada en la estructura del
algoritmo. Sus resultados se presentan más adelante como comparación complementaria
y no como reemplazo de la propuesta principal.

En conjunto, la implementación desarrollada organiza el almacenamiento comprimido
como una combinación de interfaz, backend, layout de bloques, caché y mecanismos de
ejecución de compuertas dentro de TMFQSfullstate. Esta estructura permite integrar
Blosc2 en la ruta de simulación sin modificar el modelo de estado completo y prepara la
base para la evaluación experimental presentada en el capítulo siguiente.
